% 第 13 讲：同伦余不变商、交换化与 Tate 上同调
% 源：src/lecture13.txt（235 行，含 2 处页眉）
\section*{第 13 讲\quad 同伦余不变商、交换化与 Tate 上同调}

回忆上一讲我们显式构造了 $G$-模复形 $X$ 的同伦不变量 $X^{hG}$。为此我们构造了棒消解 $P_G^{\mathrm{can}}\xrightarrow{\ \sim\ }\Z$，其中 $P_G^{\mathrm{can}}$ 是一个典范的、位于非正次数的自由 $G$-模复形。于是有拟同构
\[
X^{hG}\simeq\operatorname{Hom}_G\bigl(P_G^{\mathrm{can}},X\bigr).
\]
具体地，对 $P_G^{\mathrm{can}}$ 有
\[
\begin{array}{ccccccc}
\cdots & \longleftarrow & \Z[G^{3}] & \longleftarrow & \Z[G\times G] & \longleftarrow & \Z[G]\longleftarrow 0\longleftarrow\cdots\\[2pt]
\cdots & \longleftarrow & 0 & \longleftarrow & 0 & \longleftarrow & \Z\longleftarrow 0\longleftarrow\cdots
\end{array}
\]
两行由增广映射 $\Z[G]\to\Z$ 连接，微分为 $d_{-1}(g_1,g_2)=g_1g_2-g_1$（$G$-作用始终作用在第一个分量上）。注意若 $G$ 有限，则这些都是有限秩 $G$-模。

对每个 $G$-模 $M$，我们有
\[
\cdots\to 0\to M\xrightarrow{\ m\mapsto(gm-m)_{g\in G}\ }\underbrace{\prod_{g\in G}M}_{\{\varphi\colon G\to M\}}\to\prod_{g,h\in G}M\to\cdots,
\]
% [原文疑误：此段三处的 $M_{hG}$ 均应为 $M^{hG}$——该复形是 $\operatorname{Hom}_G(P_G^{\mathrm{can}},M)$（同伦不变子模），其零阶上同调为不变子模 $M^G$、一阶为通常的 $H^1(G,M)$；已径改]
经由某个进一步的微分，得到 $M^{hG}$。我们可以用这个表达式显式计算 $M^{hG}$ 的一阶上同调。结果表明，函数 $\varphi\colon G\to M$ 被该微分杀掉，当且仅当它是 $1$-上闭链（有时称为扭同态），即对一切 $g,h\in G$ 有
\[
\varphi(gh)=\varphi(g)+g\cdot\varphi(h)
\]
（群作用意义下）。类似地，$\varphi$ 是 $1$-上边缘，如果存在 $m\in M$ 使得对一切 $g\in G$ 有 $\varphi(g)=g\cdot m-m$。要点是
\[
H^{1}(G,M):=H^{1}(M^{hG})=\{1\text{-上闭链}\}/\{1\text{-上边缘}\}.
\]
作为推论，若 $G$ 平凡地作用在 $M$ 上，则 $H^{1}(G,M)=\operatorname{Hom}_{\mathrm{Group}}(G,M)$：因为此时 $1$-上边缘全为零，而 $1$-上闭链恰好是通常的群同态。这也说明零阶上同调恰好就是不变子模，正如上一讲所示。

现在，我们的目标（很久以前就定下的）是对任意有限群 $G$ 定义 Tate 上同调与 Tate 复形。我们希望
\[
\widehat{H}^{0}(G,M)=M^{G}/N(M)=\operatorname{Coker}\bigl(M^{G}\xrightarrow{\ N\ }M^{G}\bigr),
\]
因为它推广了局部域扩张的局部类域论的核心问题。回忆 $M_{G}=M/(g-1)M$（等价于与平凡模作张量，且与不变子模对偶——我们更愿意把后者看作子模），于是这个映射穿过 $M_{G}$，并诱导出上面的范映射。

我们的计划是：对 $G$-模复形 $X$，构造
\[
X_{hG}\xrightarrow{\ N\ }X^{hG}\longrightarrow X_{tG}:=\operatorname{hCoker}(N).
\]
因此我们首先需要定义同伦余不变商 $X_{hG}$。

注意若 $M$ 是 $G$-模，则 $M_{G}=M\otimes_{\Z[G]}\Z$。令 $I_{G}:=\ker(\varepsilon)$，于是有短正合列
\[
0\to I_{G}\to\Z[G]\xrightarrow{\ \varepsilon\ }\Z\to 0,\qquad \sum_{i}n_{i}g_{i}\mapsto\sum_{i}n_{i}.
\]
我们断言 $I_{G}$ 作为 $\Z$-模由 $\{g-1:g\in G\}$ 张成（留作习题）。一个推论是
\[
\Z[G]^{\oplus G}\to\Z[G]\to\Z\to 0
\]
正合，因为 $\Z[G]^{\oplus G}\cong I_{G}$，其中第 $g$ 个坐标上的 $1$ 映到 $g-1$。

\textbf{注 13.1.} 计算张量积的正确算法如下：回忆张量积是右正合的，即它保持满射，且与代数作张量给出原来的模。要与一个模作张量，先取该模的生成元与关系，用它写出如上的一条消解，再与这条消解作张量，得到一个定义在该模的若干拷贝的直和上的矩阵，然后取余核。

如果我能用类似的过程对链复形定义 $M_{hG}$，那将非常方便。

\textbf{定义 13.2.} 若 $X$ 与 $Y$ 是链复形，则
\[
(X\otimes Y)_{i}:=\bigoplus_{j\in\Z}X_{j}\otimes Y_{i-j},
\]
微分为
\[
d(x\otimes y):=dx\otimes y+(-1)^{j}x\otimes dy.
\]
若 $X$ 是右 $A$-模复形、$Y$ 是左 $A$-模复形，则 $X\otimes_{A}Y$ 类似地定义。

注意因子 $(-1)^{j}$ 保证微分的平方为零。另外，对代数而言不必担心左、右 $A$-模之别，因为把乘法次序反过来就给出左代数与右代数之间的同构；对 $G$-模而言，这意味着把每个元素换成它的逆。

现在，对 $X_{hG}$ 的一个糟糕猜测是 $X\otimes_{\Z[G]}\Z$，因为它不保持无环复形，等价地不保持拟同构。更好的猜测是取投射消解 $P_{G}\simeq\Z$（例如 $P_G^{\mathrm{can}}$），并用它作张量：
\[
X_{hG}:=X\otimes_{\Z[G]}P_{G}.
\]

\textbf{定义 13.3.} 左 $A$-模复形 $F$ 称为\textbf{平坦的}，如果对每个右 $A$-模的无环复形 $Y$，$Y\otimes_{A}F$ 也无环，即 $-\otimes_{A}F$ 保持单射。

我们现在问 $P_{G}$ 是否平坦。事实上：

\textbf{断言 13.4.} 任何投射复形都平坦。

一个更容易的断言是：

\textbf{断言 13.5.} 任何上有界且各项 $F^{i}$ 都平坦的复形 $F$ 都平坦。

为证明这个断言，我们要用到投射模平坦这一事实：它们是自由模的直和项，而自由模显然平坦（即若 $F=F_{1}\oplus F_{2}$，则 $F\otimes M=(F_{1}\otimes M)\oplus(F_{2}\otimes M)$）。

\textit{证明.} 情形 1. 设 $F$ 只出现在 $0$ 次，即 $i\neq 0$ 时 $F^{i}=0$。对每个复形 $Y=Y^{\bullet}$，$Y\otimes_{A}F$ 为
\[
\cdots\to Y^{i}\otimes_{A}F\xrightarrow{\ d^{i}\otimes\mathrm{id}_{F}\ }Y^{i+1}\otimes_{A}F\to\cdots.
\]
由于 $F$ 平坦，对每个 $i$ 有 $H^{i}(Y\otimes_{A}F)=H^{i}(Y)\otimes_{A}F$（因为 $F$ 平坦意味着与 $F$ 作张量与取核、余核、像交换），故若 $Y$ 无环，则 $Y\otimes_{A}F$ 也无环。

情形 2. 设 $F$ 只出现在 $0$ 与 $-1$ 次，即 $F$ 形如
\[
\cdots\to 0\to F^{-1}\to F^{0}\to 0\to\cdots,
\]
从而 $F^{\bullet}=\operatorname{hCoker}(F^{-1}\to F^{0})$。由于张量积与同伦余核交换，我们得到
\[
Y\otimes_{A}F=\operatorname{hCoker}\bigl(Y\otimes_{A}F^{-1}\to Y\otimes_{A}F^{0}\bigr),
\]
于是由情形 1，若 $Y$ 无环，则 $Y\otimes_{A}F^{0}$ 与 $Y\otimes_{A}F^{-1}$ 也无环，从而由上同调长正合列，$Y\otimes_{A}F$ 无环。类似的（归纳）论证给出 $F$ 有界的情形。

情形 3. 一般情形，考虑截断复形
\[
\begin{array}{l}
F^{[0]}=\bigl(\cdots\to 0\to 0\to F^{0}\to 0\to\cdots\bigr),\\[2pt]
F^{[1]}=\bigl(\cdots\to 0\to F^{-1}\to F^{0}\to 0\to\cdots\bigr),\\[2pt]
F^{[2]}=\bigl(\cdots\to F^{-2}\to F^{-1}\to F^{0}\to 0\to\cdots\bigr),\ \dots
\end{array}
\]
相邻截断之间的自然映射 $F^{[i]}\to F^{[i+1]}$ 在 $-i+1,\dots,0$ 次为恒等映射、在 $-i$ 次为微分 $d$；原图中的所有方块显然交换，故这些都是复形的态射，且 $F=\varinjlim_{i}F^{[i]}$。由于正向极限与张量积交换（注意逆向极限不然，因为满射性），有 $Y\otimes_{A}F=\varinjlim_{i}Y\otimes_{A}F^{[i]}$。由情形 2，每个 $Y\otimes_{A}F^{[i]}$ 无环；又上同调与正向极限交换（因为后者保持核、余核与像），故若 $Y$ 无环，则 $Y\otimes_{A}F$ 也无环。\hfill$\square$

\textbf{注 13.6.} 设 $Y$ 是 $A$-模复形，取拟同构 $F\xrightarrow{\ \sim\ }Y$，其中 $F$ 平坦，并定义
\[
Y\otimes_{A}^{\mathrm{der}}X:=F\otimes_{A}X.
\]
则它在拟同构意义下良定义，而拟同构又在同伦意义下良定义，如此以往（真是一路到底！）。

\textbf{定义 13.7.}（$Y$ 相对于 $X$ 的）第 $i$ 个\textbf{挠群}为
\[
\operatorname{Tor}_{i}^{A}(Y,X):=H^{-i}\bigl(Y\otimes_{A}^{\mathrm{der}}X\bigr).
\]

\textbf{定义 13.8.} 链复形 $X$ 的\textbf{同伦余不变商}是复形
\[
X_{hG}:=X\otimes_{\Z[G]}^{\mathrm{der}}\Z\simeq X\otimes_{\Z[G]}P_{G}
\]
（注意它只在拟同构意义下良定义）。

\textbf{定义 13.9.} $H_{i}(G,X):=H^{-i}(X_{hG})$（注意这里更偏好下标记号，因为 $X_{hG}$ 一般只在非正次数上有定义）。

下面做一些基本计算。

\textbf{断言 13.10.} 若 $X$ 被 $0$ 从上界住，则 $H_{0}(G,X)=H_{0}(X)_{G}$（证明与断言 12.5 类似）。

\textbf{断言 13.11.} $H_{1}(G,\Z)=G^{\mathrm{ab}}$，其中 $G^{\mathrm{ab}}$ 表示 $G$ 的交换化。

注意这可以说是与本讲开头所见结论相对偶的陈述；$H^{1}(G,M)$ 关乎映射 $G\to M$，这等同于映射 $G^{\mathrm{ab}}\to M$，而这里 $H_{1}(G,\Z)$ 由走出 $G$ 的映射决定。

\textit{证明.} 回忆短正合列
\[
0\to I_{G}\to\Z[G]\xrightarrow{\ \varepsilon\ }\Z\to 0.
\]
上同调长正合列给出正合列
\[
H_{1}(G,\Z[G])\to H_{1}(G,\Z)\to H_{0}(G,I_{G})\to H_{0}(G,\Z[G])\to H_{0}(G,\Z).
\]
我们有
\[
H_{0}(G,\Z[G])=H_{0}(\Z[G])_{G}=\Z[G]\otimes_{\Z[G]}\Z=\Z
\]
（由断言 13.10）。当然 $H_{0}(G,\Z)=H_{0}(\Z)_{G}=\Z$，且 $H_{1}(G,\Z[G])=0$，因为
\[
\Z[G]_{hG}:=\Z[G]\otimes_{\Z[G]}P_{G}=P_{G}\simeq\Z
\]
是拟同构。于是我们的正合列实际上是
\[
0\to H_{1}(G,\Z)\xrightarrow{\ \sim\ }H_{0}(G,I_{G})\to\Z\xrightarrow{\ \sim\ }\Z,
\]
这给出所说的同构。要点是
\[
H_{1}(G,\Z)=(I_{G})_{G}=I_{G}/I_{G}^{2},
\]
因为 $M_{G}=M/I_{G}\cdot M$。

\textbf{断言 13.12.} 映射
\[
\Z[G]/I_{G}^{2}\to G^{\mathrm{ab}}\times\Z,\qquad g\mapsto(\bar g,1)
\]
是同构。

这将推出 $I_{G}/I_{G}^{2}=\ker(\varepsilon)/I_{G}^{2}=G^{\mathrm{ab}}$，正是所要。

\textit{证明.} 首先注意上述映射是同态。事实上，记 $[g]\in\Z[G]$ 为 $g$ 的类，则对任意 $g,h\in G$ 有
\[
[g]+[h]\mapsto(\bar g\bar h,2),\qquad [g]\mapsto(\bar g,1),\qquad [h]\mapsto(\bar h,1),
\]
而后两个像之和恰为第一个。我们断言该映射有逆，由
\[
G\times\Z\to\Z[G]/I_{G}^{2},\qquad (g,n)\mapsto[g]+n-1
\]
诱导。它是同态，因为
\[
([g]-1)([h]-1)=[gh]-[g]-[h]+1\in I_{G}^{2},
\]
从而
\[
([g]-1)+([h]-1)\equiv[gh]-1\pmod{I_{G}^{2}},
\]
正是所要。最后，二者互逆，因为
\[
(\bar g,1)\mapsto[g]+1-1=[g],\qquad [g]+n-1\mapsto(\bar g,1)(1,n-1)=(\bar g,n),
\]
正是所要。\hfill$\square$

这证明了该断言。

最后，我们定义范映射 $X_{hG}\xrightarrow{\ N\ }X^{hG}$ 为如下复合：
\[
X_{hG}=X\otimes_{\Z[G]}P_{G}\to X\otimes_{\Z[G]}\Z\to\operatorname{Hom}_{\Z[G]}(\Z,X)\to\operatorname{Hom}_{\Z[G]}(P_{G},X)=X^{hG},
\]
其中第二个映射经由逐项的范映射（用张量–同调伴随）。然后置
\[
X_{tG}:=\operatorname{hCoker}\bigl(X_{hG}\xrightarrow{\ N\ }X^{hG}\bigr),
\]
我们断言它在拟同构意义下推广了先前对循环群所得者，于是可以定义
\[
\widehat{H}^{i}(G,X):=H^{i}(X_{tG}).
\]

很快我们将证明：

\textbf{断言 13.13（lcft）.} 对有限群 $G$ 与局部域扩张 $L/K$，$P_{G}\to L^{\times}[2]$ 在 Tate 上同调上是同构。

这给出
\[
\widehat{H}^{-2}(G,\Z)\cong\widehat{H}^{0}(G,L^{\times})=K^{\times}/N(L^{\times}).
\]
我们有正合列
\[
0=H^{-2}(\Z_{hG})\to\widehat{H}^{-2}(G,\Z)\xrightarrow{\ \sim\ }\underbrace{H^{-1}(\Z_{hG})}_{H_{1}(G,\Z)=G^{\mathrm{ab}}}\to H^{-1}(\Z^{hG})=0,
\]
因为 $\Z_{hG}$ 位于非正次数。于是，对 Galois 群为 $G$ 的局部域扩张 $L/K$，有
\[
K^{\times}/N(L^{\times})\cong G^{\mathrm{ab}}.
\]
